% 统计图重绘：2016 北京文科卷第 17 题频率分布直方图。
% 真值来源：source/普通高考/2016/2016北京文.pdf，第 1 页；pdfimages -list 未发现嵌入图像。
% 定位证据：tmp/redraw_2016_beijing_liberal/page-1-grid100.jpg；无网格裁切证据：
% tmp/redraw_2016_beijing_liberal/q17-finalcrop.png。图中横轴为用水量（立方米），
% 纵轴为频率/组距；组距均为 0.5，八个区间柱高依次为 0.2、0.3、0.4、0.5、
% 0.3、0.1、0.1、0.1。水平虚线分别指向对应柱顶。
\begin{tikzpicture}
  \begin{axis}[exam statistical histogram,
    width=12.7cm,
    height=7.4cm,
    xmin=0, xmax=5.35,
    ymin=0, ymax=0.58,
    axis lines=left,
    axis line style={-stealth, line width=0.7pt},
    clip=false,
    xtick={0,0.5,1,1.5,2,2.5,3,3.5,4,4.5},
    xticklabels={0,0.5,1,1.5,2,2.5,3,3.5,4,4.5},
    ytick={0.1,0.2,0.3,0.4,0.5},
    yticklabels={0.1,0.2,0.3,0.4,0.5},
    tick style={draw=none},
    tick label style={font=\normalsize},
  ]
    % 原图中的水平虚线基准。
    \draw[dashed, line width=0.55pt] (axis cs:0,0.1) -- (axis cs:4.5,0.1);
    \draw[dashed, line width=0.55pt] (axis cs:0,0.2) -- (axis cs:1,0.2);
    \draw[dashed, line width=0.55pt] (axis cs:0,0.3) -- (axis cs:3,0.3);
    \draw[dashed, line width=0.55pt] (axis cs:0,0.4) -- (axis cs:2,0.4);
    \draw[dashed, line width=0.55pt] (axis cs:0,0.5) -- (axis cs:2.5,0.5);
    % 八个等组距矩形柱。
    \addplot[/tikz/ybar interval,mark=none,draw=black,fill=none] coordinates {
      (0.5,0.2)
      (1,0.3)
      (1.5,0.4)
      (2,0.5)
      (2.5,0.3)
      (3,0.1)
      (3.5,0.1)
      (4,0.1)
      (4.5,0)
    };
    \examstatylabel{\(\dfrac{\text{频率}}{\text{组距}}\)};
    \examstatxlabel{用水量（立方米）};
  \end{axis}
\end{tikzpicture}
